\subsection{Mejoras de resultados anteriores.}
\label{subsection:IV.1.7-fr}

Vamos a dar en este número una lista de proposiciones demostradas en
los capítulos precedentes, cuyo enunciado puede mejorarse mediante las
nuevas condiciones de finitud introducidas anteriormente.

\begin{env}[1.7.1]
\label{IV.1.7.1-fr}
En los enunciados de (\textbf{I}, 6.4.2, 6.4.3 y 6.4.9), se puede, en lugar de suponer
que $X$ y $Y$ son de tipo finito sobre el cuerpo $K$, suponer solamente que
son \emph{localmente de tipo finito} sobre $K$; para (\textbf{I}, 6.4.2), basta
observar que todo punto de un preesquema $X$ que es cerrado en todo abierto
afín de $X$ es cerrado en $X$, y un esquema afín localmente de tipo finito
sobre $K$ es \emph{ipso facto} de tipo finito (1.5.1). Para (\textbf{I}, 6.4.9), se utiliza
(1.3.7) y (1.3.4, (v)).
\end{env}

\begin{env}[1.7.2]
\label{IV.1.7.2-fr}
En (\textbf{I}, 6.5.1, (ii)), se puede, en lugar de suponer que $S$ es localmente
noetheriano y $Y$ de tipo finito sobre $S$, suponer solamente que $Y$ es
\emph{localmente de presentación finita} sobre $S$, como muestra inmediatamente la
demostración (aplicando la definición (1.4.2)). Asimismo, en (\textbf{I}, 6.5.4,
(ii)) y (\textbf{I}, 6.5.5, (ii)) basta suponer que $f$ es un $S$-morfismo,
siendo $Y$ \emph{localmente de presentación finita} sobre $S$.
\end{env}

\begin{env}[1.7.3]
\label{IV.1.7.3-fr}
En (\textbf{I}, 7.1.11), para la segunda afirmación, en lugar de suponer $S$
localmente noetheriano y $Y$ de tipo finito sobre $S$, se puede suponer solamente
$Y$ localmente de presentación finita sobre $S$ (pues la demostración depende de
(\textbf{I}, 6.5.1, (ii))); asimismo en (\textbf{I}, 7.1.12 a 7.1.15).
\end{env}

\begin{env}[1.7.4]
\EGAReaderTarget{ega4.c04.statement.1-7-4}{IV.1.7.4-fr}
En (\textbf{I}, 9.2.2), se pueden sustituir las tres hipótesis por la sola hipótesis
(que se deduce de cada una de las otras) de que $f$ es \emph{cuasicompacto y
cuasiseparado}, en virtud de (1.2.6) y (1.2.7). Se observará que en (\textbf{I}, 9.2.1),
la hipótesis significa exactamente que $f$ es cuasicompacto y cuasiseparado.
\end{env}

\begin{env}[1.7.5]
\label{IV.1.7.5-fr}
En (\textbf{I}, 9.3.1, 9.3.2 y 9.3.3), se pueden sustituir las dos hipótesis sobre $X$
por la hipótesis (más débil) de que $X$ es un preesquema \emph{cuasicompacto y
cuasiseparado}, pues la demostración utiliza exactamente estas dos propiedades (en
virtud de (1.2.7)).
\end{env}

\begin{env}[1.7.6]
\label{IV.1.7.6-fr}
En (\textbf{I}, 9.3.4 y 9.3.5), basta suponer $X$ cuasicompacto,
$\mathcal F$ y $\mathcal J$ cuasicoherentes y de tipo finito.
\end{env}

\begin{env}[1.7.7]
\label{IV.1.7.7-fr}
En (\textbf{I}, 9.4.7), se puede debilitar la hipótesis \emph{b)} suponiendo solamente $X$
\emph{cuasiseparado}, puesto que solo se necesita que la inmersión canónica
$U\to X$ sea cuasicompacta (1.2.7). Se puede hacer la misma modificación en
(\textbf{I}, 9.4.8, 9.4.9 y 9.4.10).
\end{env}

\begin{env}[1.7.8]
\label{IV.1.7.8-fr}
En (\textbf{I}, 9.5.1), las condiciones indicadas entre paréntesis en el enunciado
pueden sustituirse por la condición de que $f$ sea cuasicompacto y
cuasiseparado.
\end{env}

\begin{env}[1.7.9]
\label{IV.1.7.9-fr}
En (\textbf{I}, 9.6.5), se pueden sustituir las hipótesis sobre $X$ por la hipótesis
más débil de que $X$ es \emph{cuasicompacto y cuasiseparado}, como muestra la
demostración,
\oldpage[IV-1]{237}
que solo utiliza (\textbf{I}, 9.4.7). En (\textbf{I}, 9.6.6), la hipótesis «$X$ es un
esquema cuasicompacto» puede sustituirse por «$X$ es un preesquema
cuasicompacto y cuasiseparado».
\end{env}

\begin{env}[1.7.10]
\label{IV.1.7.10-fr}
En (\textbf{II}, 1.7.15), se puede suponer solamente que $Y$ es cuasicompacto y
cuasiseparado, pues la demostración solo utiliza (\textbf{I}, 9.6.5).
\end{env}

\begin{env}[1.7.11]
\label{IV.1.7.11-fr}
En la segunda afirmación de (\textbf{II}, 3.4.5), se pueden sustituir las dos
hipótesis sobre $Y$ por la hipótesis más débil de que $Y$ es cuasicompacto y
cuasiseparado. Asimismo en (\textbf{II}, 3.4.8).
\end{env}

\begin{env}[1.7.12]
\label{IV.1.7.12-fr}
En (\textbf{II}, 3.8.5), se puede igualmente suponer solamente que $Y$ es cuasicompacto
y cuasiseparado, pues esta hipótesis interviene a través de (\textbf{I}, 9.4.7).
\end{env}

\begin{env}[1.7.13]
\label{IV.1.7.13-fr}
En (\textbf{II}, 4.4.1, 4.4.6 y 4.4.7), basta suponer que $Y$ es
\emph{cuasicompacto y cuasiseparado}, teniendo en cuenta (1.2.3) en la
demostración de (\textbf{II}, 4.4.7). Asimismo, en (\textbf{II}, 4.4.10.1), basta
suponer $Z$ cuasicompacto y cuasiseparado.
\end{env}

\begin{env}[1.7.14]
\label{IV.1.7.14-fr}
En (\textbf{II}, 4.5.2) y (\textbf{II}, 4.5.5), basta suponer $X$
\emph{cuasicompacto y cuasiseparado}, pues esta hipótesis interviene a través de (\textbf{I},
9.3.1 y 9.3.2).
\end{env}

\begin{env}[1.7.15]
\label{IV.1.7.15-fr}
En (\textbf{II}, 4.6.8), basta de nuevo suponer $X$ \emph{cuasicompacto y
cuasiseparado}, pues esta hipótesis interviene a través de (\textbf{I}, 9.4.7).
\end{env}

\begin{env}[1.7.16]
\label{IV.1.7.16-fr}
En (\textbf{II}, 5.1.2), basta suponer que $X$ es cuasicompacto y
cuasiseparado (la hipótesis interviene a través de (\textbf{II}, 4.5.2 y 4.5.5)). Asimismo en
(\textbf{II}, 5.1.9), donde se utiliza (\textbf{I}, 9.6.5).
\end{env}

\begin{env}[1.7.17]
\label{IV.1.7.17-fr}
En (\textbf{II}, 5.2.1), sigue bastando suponer $X$ cuasicompacto y
cuasiseparado (la hipótesis interviene a través de (\textbf{II}, 4.5.2)).
\end{env}

\begin{env}[1.7.18]
\label{IV.1.7.18-fr}
En (\textbf{II}, 5.2.2), hay que suponer $f$ \emph{cuasicompacto} y $X$ \emph{cuasiseparado}; la
demostración actual es, en efecto, insuficiente (con las hipótesis formuladas),
pues del hecho de que para todo $\mathcal O_X$-Módulo cuasicoherente $\mathscr F$ se
tenga $\mathrm R^1f_*(\mathscr F)=0$, no se sigue necesariamente, para un
abierto afín $U$ de $Y$, que la misma propiedad sea válida para los
$(\mathcal O_X|f^{-1}(U))$-Módulos cuasicoherentes, si no se sabe que uno de tales
Módulos es la \emph{restricción} de un $\mathcal O_X$-Módulo cuasicoherente (la
misma observación se aplica a las condiciones \emph{b)} y \emph{c')}); cuando $f$ es
cuasicompacto y $X$ cuasiseparado, esta prolongación es posible en virtud de
(\textbf{I}, 9.4.7), modificado anteriormente (1.7.7).
\end{env}

\begin{env}[1.7.19]
\label{IV.1.7.19-fr}
En (\textbf{II}, 5.3.2), basta suponer $Y$ \emph{cuasicompacto y
cuasiseparado} (lo que interviene a través de (\textbf{II}, 4.4.6 y 4.4.7)). Asimismo en (\textbf{II},
5.5.3, (ii)).
\end{env}

\begin{env}[1.7.20]
\label{IV.1.7.20-fr}
En (\textbf{III}, 1.2.2, 1.2.3 y 1.2.4), se pueden sustituir las dos hipótesis sobre
$X$ por la hipótesis más débil de que $X$ es un preesquema \emph{cuasicompacto
y cuasiseparado}, pues la demostración solo utiliza (\textbf{I}, 9.3.3).
\end{env}

\begin{env}[1.7.21]
\label{IV.1.7.21-fr}
En (\textbf{III}, 1.4.10 a 1.4.14), se puede sustituir «separado» por
«cuasiseparado», pues esta hipótesis sirve únicamente para poder aplicar (\textbf{I},
9.2.2) (véase anteriormente (1.7.4)). Asimismo, en (\textbf{III}, 1.4.15), basta
suponer que el morfismo $u$ es \emph{cuasiseparado y cuasicompacto}.
\end{env}

\begin{env}[1.7.22]
\label{IV.1.7.22-fr}
En (\textbf{III}, 2.1.3), se pueden de nuevo sustituir las hipótesis por la hipótesis
más débil de que $X$ es un preesquema cuasicompacto y cuasiseparado, siendo el
razonamiento el mismo que en (\textbf{III}, 1.2.2).
\end{env}

\oldpage[IV-1]{238}
